%%%%%%%%%%%%%%%%
%%% exod 11 %%%
%%%%%%%%%%%%%%%%

\Chapter{11}

\Verse{1}
{
	\hRJE{וַיֹּאמֶר יְהֹוָה אֶל מֹשֶׁה עוֹד נֶגַע אֶחָד אָבִיא עַל פַּרְעֹה וְעַל מִצְרַיִם אַחֲרֵי כֵן יְשַׁלַּח אֶתְכֶם מִזֶּה כְּשַׁלְּחוֹ כָּלָה גָּרֵשׁ יְגָרֵשׁ אֶתְכֶם מִזֶּה׃}
}
{
	\eRJE{
		And YHWH said to Moses, “I’ll bring one more plague on
		Pharaoh and on Egypt. After that he’ll let you go from
		here. When he lets go completely he’ll drive you out from
		here!
	}
}
{
}

\Verse{2}
{
	\hRJE{דַּבֶּר נָא בְּאזְנֵי הָעָם וְיִשְׁאֲלוּ אִישׁ מֵאֵת רֵעֵהוּ וְאִשָּׁה מֵאֵת רְעוּתָהּ כְּלֵי כֶסֶף וּכְלֵי זָהָב׃}
}
{
	\eRJE{
		Speak in the people’s ears that each man will ask of his
		neighbor and each woman of her neighbor items of silver
		and items of gold.”
	}
}
{
}

\Verse{3}
{
	\hRJE{וַיִּתֵּן יְהֹוָה אֶת חֵן הָעָם בְּעֵינֵי מִצְרָיִם גַּם הָאִישׁ מֹשֶׁה גָּדוֹל מְאֹד בְּאֶרֶץ מִצְרַיִם בְּעֵינֵי עַבְדֵי פַרְעֹה וּבְעֵינֵי הָעָם׃}
}
{
	\eRJE{
		And YHWH put the people’s favor in the Egyptians’ eyes.
		Also, the man Moses was very big in the land of Egypt in
		the eyes of Pharaoh’s servants and in the people’s eyes.
	}
}
{
}

\Verse{4}
{
	\hRJE{וַיֹּאמֶר מֹשֶׁה כֹּה אָמַר יְהֹוָה כַּחֲצֹת הַלַּיְלָה אֲנִי יוֹצֵא בְּתוֹךְ מִצְרָיִם׃}
}
{
	\eRJE{
		And Moses said, “YHWH said this: In the middle of the
		night I am going out through Egypt,
	}
}
{
%	\footnote{
%		I am going out through Egypt. YHWH Himself. Later Jewish tradition,
%		presumably unable to bear the thought of God personally passing through
%		Egypt and causing the deaths, introduced the horrible concept of the
%		“Angel of Death.” But there is no such thing as an Angel of Death in the
%		Bible. The text is explicit that God personally passes through Egypt.
%		The subtlety comes in the question of the relationship between the deity
%		and the agent of the destruction: on one hand, God says, “I strike
%		Egypt” (12:13), and, on the other hand, it appears that God halts at the
%		door and prevents this destroyer (mašît) from coming into the house. Is
%		the mašît a sickness? a force? or a personified being? The wording of
%		12:13 (“a plague among you as a destroyer”) indicates that the destroyer
%		is the plague itself and not some person or angel. In any case, God
%		clearly is responsible for the creation of it and can therefore say “I
%		strike.”
%	}
}

\Verse{5}
{
	\hRJE{וּמֵת כּל בְּכוֹר בְּאֶרֶץ מִצְרַיִם מִבְּכוֹר פַּרְעֹה הַיֹּשֵׁב עַל כִּסְאוֹ עַד בְּכוֹר הַשִּׁפְחָה אֲשֶׁר אַחַר הָרֵחָיִם וְכֹל בְּכוֹר בְּהֵמָה׃}
}
{
	\eRJE{
		and every firstborn in the land of Egypt will die, from
		the firstborn of Pharaoh who is sitting on his throne to
		the firstborn of the maid who is behind the mill and every
		firstborn of an animal.
	}
}
{
}

\Verse{6}
{
	\hRJE{וְהָיְתָה צְעָקָה גְדֹלָה בְּכל אֶרֶץ מִצְרָיִם אֲשֶׁר כָּמֹהוּ לֹא נִהְיָתָה וְכָמֹהוּ לֹא תֹסִף׃}
}
{
	\eRJE{
		And there will be a big cry in all the land of Egypt, that
		there has been none like it and there won’t continue to be
		like it.
	}
}
{
}

\Verse{7}
{
	\hRJE{וּלְכֹל בְּנֵי יִשְׂרָאֵל לֹא יֶחֱרַץ כֶּלֶב לְשֹׁנוֹ לְמֵאִישׁ וְעַד בְּהֵמָה לְמַעַן תֵּדְעוּן אֲשֶׁר יַפְלֶה יְהֹוָה בֵּין מִצְרַיִם וּבֵין יִשְׂרָאֵל׃}
}
{
	\eRJE{
		But not a dog will move its tongue at any of the children
		of Israel, from man to animal, so that you’ll know that
		YHWH will distinguish between Egypt and Israel.
	}
}
{
}

\Verse{8}
{
	\hRJE{וְיָרְדוּ כל עֲבָדֶיךָ אֵלֶּה אֵלַי וְהִשְׁתַּחֲווּ לִי לֵאמֹר צֵא אַתָּה וְכל הָעָם אֲשֶׁר בְּרַגְלֶיךָ וְאַחֲרֵי כֵן אֵצֵא וַיֵּצֵא מֵעִם פַּרְעֹה בּחֳרִי אָף׃}
}
{
	\eRJE{
		And all these servants of yours will come down to me, and
		they’ll bow to me, saying: ‘Go out, you and all the people
		who are at your feet.’ And after that I’ll go out!” And he
		went out from Pharaoh in a flaring of anger.
	}
}
{
%	\footnote{
%		they’ll bow to me. Read this speech carefully. At the beginning, when
%		Moses uses the first person he is quoting God: “I am going out through
%		Egypt . . .” But here at the end, the first person refers to himself:
%		“will come down to me, and they’ll bow to me. . . . And after that I’ll
%		go out!” It is hard to say where in the text he stops quoting God and
%		starts referring to himself. This is tremendously important, a turning
%		point in Moses’ life. His fear of Pharaoh and his lack of confidence are
%		gone. And he has shifted Pharaoh’s attention from God to himself. He has
%		not stepped over the line and failed to credit God, but he will do that
%		years later in his sin at Meribah (Numbers 20). This moment in the
%		Pharaoh’s court is thus a step in the long process of Moses’ evolving
%		relationship with his God and his people, which will have a tragic
%		aspect in the end.
%	}
}

\Verse{9}
{
	\hRJE{וַיֹּאמֶר יְהֹוָה אֶל מֹשֶׁה לֹא יִשְׁמַע אֲלֵיכֶם פַּרְעֹה לְמַעַן רְבוֹת מוֹפְתַי בְּאֶרֶץ מִצְרָיִם׃}
}
{
	\eRJE{
		And YHWH said to Moses, “Pharaoh won’t listen to you—in
		order to multiply my wonders in the land of Egypt.”
	}
}
{
}

\Verse{10}
{
	\hRJE{וּמֹשֶׁה וְאַהֲרֹן עָשׂוּ אֶת כּל הַמֹּפְתִים הָאֵלֶּה לִפְנֵי פַרְעֹה וַיְחַזֵּק יְהֹוָה אֶת לֵב פַּרְעֹה וְלֹא שִׁלַּח אֶת בְּנֵי יִשְׂרָאֵל מֵאַרְצוֹ׃}
}
{
	\eRJE{
		And Moses and Aaron had done all these wonders in front of
		Pharaoh, and YHWH strengthened Pharaoh’s heart, and he did
		not let the children of Israel go from his land.
	}
}
{
}
